%% ===========================================================================
\FloatBarrier
\section{Numerical verification and limiting cases}
\label{sec:verification}
%% ===========================================================================

\FloatBarrier
\subsection{Material, geometry and discretization}
\label{sec:material}

The four verification cases use the same baseline material and geometry,
with only the parameter under examination varied. The baseline material is
room-temperature tetragonal $\mathrm{BaTiO_3}$.

The dielectric anisotropy ratio
$\kappa_{11}/\kappa_{33}=39$ is retained, so the two poling configurations
differ in their permittivity as well as in their modulus. The strain-gradient
length scale is set by the window: with half-widths
$a_1\times a_2=3\times1$\,\si{\micro\meter} the cuboid window of
\eqref{eq:Mshapes} gives $M_{11}=3.0\times10^{-12}$ and
$M_{22}=3.3\times10^{-13}$\,\si{\square\meter}, that is characteristic lengths
$\sqrt{M_{11}}=\SI{1.73}{\micro\meter}$ and
$\sqrt{M_{22}}=\SI{0.58}{\micro\meter}$ along and across the strip. The ratio
$a_1/L=0.015$ lies within the assumed input range $a/L\le0.2$.
At fixed stiffness and geometry, the angular frequencies scale as
$\omega\propto\rho^{-1/2}$ when the density changes uniformly, with
$f=\omega/(2\pi)$; this common density scaling alone does not change the
open-/short-circuit coupling factors. The flexoelectric coefficient $\mu$ enters the coupling matrix linearly.
For an isolated flexoelectric mechanism its small-coupling response has a
leading $\mu^2$ contribution, but the coupled model also contains
piezoelectric--flexoelectric cross terms. The reported growth factors and
comparisons are therefore conditional on the assumed $\mu$, rather than
independent of its value.

\FloatBarrier
\subsection{Case I: benchmarks and mesh convergence}

Case I compares finite-element calculations with available analytical
benchmarks and assesses mesh-refinement sensitivity. Supporting calculations
are given in \ref{app:verification}.
In every case the solver assembles the weak form \eqref{eq:weak} with the
Hermite interpolation \eqref{eq:interp}, applies the kinematic conditions
\eqref{eq:bcmech}, and returns the higher-order stress \eqref{eq:tau} from the
constitutive relations of \eqref{eq:govexp}.

Three problems with known answers were solved with the same element: a
one-dimensional strain-gradient bar, whose static tip displacement and
frequency were compared with the closed form; the same bar in the classical
limit; and the constant-strain patch test. The clamped-free and clamped-clamped
frequencies of the strip agree with the Euler--Bernoulli values to $0.56\,\%$
and $0.46\,\%$, and the elastic strip recovers classical elasticity when
the window is set to zero. Figure~\ref{fig:convergence} then shows the error in
the frequency of the extensional mode and in its quasi-static coupling factor
against the number of degrees of freedom. Both approach their fine-mesh reference values; the plotted
$N^{-2}$ line is a comparison guide, not a fitted convergence exponent. The $48\times4$ mesh lies within the tested refinement range: refining it to $96\times4$
changes the frequency by $0.05\,\%$ and the coupling factor by $0.03\,\%$,
while the coarser $24\times2$ mesh used for the parameter sweeps differs from
it by $0.04\,\%$ and $0.03\,\%$. Table~\ref{tab:convergence} lists the four
meshes.

\begin{figure}[!htbp]
\centering
\includegraphics[width=\textwidth]{F2_convergence.pdf}
\caption{Variations of (a) the relative error of the extensional coupling factor $k^2$ and of the frequency of the mode that carries it, both referred to the finest mesh, and of (b) the two quantities themselves, against the number of mechanical degrees of freedom. The line in (a) is an $N^{-2}$ reference guide (not an independently fitted rate), the dotted line in (b) the finest-mesh reference value, and the $48\times4$ mesh is marked. Its zero self-reference error is omitted
from the logarithmic error panel. The strip uses the baseline material and geometry, poling along $x$, plane stress and rectangular window half-widths
$a_1=\SI{3}{\micro\meter}$ and $a_2=\SI{1}{\micro\meter}$.}
\label{fig:convergence}
\end{figure}

\subsection{Case II: the coupling factor and its closed form}

Case II compares the finite-element coupling factor with its
closed-form constitutive counterpart.

For the extensional electrode configuration, the constitutive coupling
factor can be calculated independently of the finite-element mesh. From the constitutive tensors alone,
\begin{equation}
  x=\frac{e^{\mathsf T}C_{p}^{-1}e}{\kappa}=1.3542,
  \qquad
  k^{2}=\frac{x}{1+x}=0.5752 ,
  \label{eq:xclosed}
\end{equation}
which is the closed-form value for the plane-stress strip; the
three-dimensional free bar gives $0.5771$, the three-dimensional uniaxial case
$0.3748$ and the two-dimensional plane-strain case $0.5132$, all from the same
material tensors under the reductions in Sec.~\ref{sec:reduction}. The plane-stress value is
the relevant one, because the strip is free out of plane.

The assembled calculation gives different values depending on the
mechanical constraints and gradient stiffness. With the window switched off and
the strip unconstrained, the assembled route reproduces \eqref{eq:xclosed} to
$9.2\times10^{-12}$ relative, at $0.57523$. This agreement checks the assembled coupling algebraically. In this free-body
calculation the inverse denotes the compatible least-squares pseudoinverse
used in the implementation (relative singular-value cutoff $10^{-11}$),
not an inverse on the rigid-body nullspace.
Clamping the left edge lowers the coupling to $0.57410$,
because the clamp restrains the in-plane deformation that carries the
piezoelectric charge, and switching the window on lowers it further to
$0.57268$, because the gradient stiffness adds energy to the same mode without
adding charge. The coupling quoted in the abstract is $k^{2}=0.5727$, which is approximately $0.44\,\%$ below the free-strip closed form and
corresponds to the specified clamped strip. The two mode-based values,
$k^2_4=0.4847$ on the elastic memory matrix and $k^2_4=0.4784$ on the
short-circuit stiffness, are different measures used in the ring-down calculations; they are not comparable
with the mode-free ones and are never averaged with them.

\subsection{Case III: time discretization and the step used}

Case III examines temporal refinement of the L1 integration and specifies
the time step used for the subsequent ring-down calculations.

The L1 discretization is assessed by step refinement: its observed order
depends on solution regularity and the measured quantity, while the
$\alpha=1$ backward-difference term is first order. At
$\alpha=1$ and $\omega\tau=0.3$ the damping ratio
extracted from the ring-down is $0.150111$, $0.150057$ and $0.150029$ at $N$,
$2N$ and $4N$ steps over twenty periods ($N=12000$, that is $600$ steps per
period), against the exact $\omega\tau/2=0.150000$: halving the step halves the
residual, an observed order of $0.97$, and the sequence extrapolates to the
exact value. The detailed ring-downs are run at the $4N$ step.
The eight transient curves use a coarser march of $400$
steps per period, giving a time-step-dependent uncertainty in the last reported decimal.

\subsection{Case IV: the fields of the coupled bending mode}

Case IV examines the displacement, potential and strain-gradient fields
of a flexoelectrically coupled bending mode.

Figure~\ref{fig:fields} records the four fields of the coupled bending mode in
the open-circuit configuration, together with the strain gradient that drives
the flexoelectric term. The in-plane and the transverse displacement are those
of a first bending mode; the potential is not. It reaches tens of millivolts a
few tens of micrometers from the clamped edge and near mid-height, and falls
toward both the edge and the free end, so the two face electrodes sit where the
potential is smallest. The field distribution is consistent with cancellation of the
net piezoelectric charge and a nonzero flexoelectric contribution, and the gradient panel shows where the flexoelectric
charge is generated: in the boundary layer at the clamp and along the neutral
axis, not at the electrodes.

\begin{figure}[!htbp]
\centering
\includegraphics[width=\textwidth]{F3_fields.pdf}
\caption{Variations of (a) the in-plane displacement, (b) the deflection, (c) the potential at \SI{10}{\nano\meter} of tip motion, (d) the electric field and (e) the magnitude of the strain gradient that drives the flexoelectric term over the coupled bending mode, in the open-circuit configuration on a $48\times4$ mesh, using the baseline geometry, material and window, poling along $y$ and plane stress. The potential is smallest at the face electrodes and the gradient largest at the clamp.}
\label{fig:fields}
\end{figure}

\FloatBarrier
\subsection{The worked calculation of the headline numbers}
\label{sec:worked}

The principal numerical results are obtained from the following
material and discretization parameters; Table~\ref{tab:worked} lists inputs and
outputs side by side.

\paragraph{Window} For half-widths $a_1=\SI{3}{\micro\meter}$ and
$a_2=\SI{1}{\micro\meter}$ the second-moment tensor of the cuboid window is
diagonal with
\begin{equation}
  M_{11}=\frac{a_1^{2}}{3}=3.000\times10^{-12}\ \si{\meter\squared},\qquad
  M_{22}=\frac{a_2^{2}}{3}=3.333\times10^{-13}\ \si{\meter\squared},
  \label{eq:Mnumber}
\end{equation}
an anisotropy ratio $M_{11}/M_{22}=9$. The three shapes are interchangeable
once $M$ is fixed, because $M_{ii}=a_i^{2}/3$, $a_i^{2}/4$ and $a_i^{2}/5$ for
the cuboid, the ellipse and the ellipsoid: an ellipsoid of half-width
$a_{\rm ell}=1.2910\,a_{\rm cub}=\sqrt{5/3}\,a_{\rm cub}$ reproduces
\eqref{eq:Mnumber} exactly, and the reciprocal
$a_{\rm cub}=0.7746\,a_{\rm ell}$ is the conversion used for the equal-$M$
window comparison.

\paragraph{Plane-stress reduction and closed form} The strip is free out of
plane, so the stiffness that multiplies $\varepsilon_{xx}$ is the reduced entry
\eqref{eq:Cp}, which for $\mathrm{BaTiO_3}$ gives
$C^{p}_{xx}=9.5500\times10^{10}$\,Pa with the poling along $x$ and
$1.6946\times10^{11}$\,Pa with the poling along $y$, and
$C^{p}_{12}=5.7000\times10^{10}$\,Pa in both cases. The unreduced 3-D record contains $e_{33}=6.7$ and $e_{31}=-0.7$\,\si{\coulomb\per\square\meter}; after the coupled plane-stress elimination, the active components are
$e^{p}=(7.05,-0.35946)$\,\si{\coulomb\per\square\meter}. The effective permittivity across the $x$-directed electrode gap is
$\kappa^{p}_{xx}=4.9804\times10^{-10}$\,\si{\farad\per\meter}. Hence
$x=1.3542$ and $k^{2}=0.5752$ from \eqref{eq:xclosed}, with no mesh involved.

\paragraph{Assembled route} The finite-element coupling is calculated by short- and open-circuit
condensation on a $48\times4$ mesh:
$245$ nodes, eight mechanical degrees of freedom per node ($1960$ unknowns)
and $245$ potential unknowns, condensed to a single electrode pair $(d,\beta)$
with $\beta=2.4902\times10^{-11}$\,\si{\farad\per\meter} in the two-dimensional, per-unit-width convention. Then
\begin{equation}
  x_{\rm FEM}=\frac{d^{\mathsf T}K_{sc}^{-1}d}{\beta}=1.3402,
  \qquad
  k^{2}_{\rm FEM}=\frac{x_{\rm FEM}}{1+x_{\rm FEM}}=0.5727 ,
  \label{eq:xFEM}
\end{equation}
the coupling factor $k^2_{\rm FEM}$ is $0.44\,\%$ below \eqref{eq:xclosed}; the clamped edge accounts for about $0.20\,\%$ of the coupling reduction
and the gradient stiffness for a further $0.25\,\%$.

\paragraph{Flexoelectric smallness and L1 weights} With the piezoelectric
constants switched off, a prescribed edge-axis field
$\varepsilon_{xx}=\kappa_c y$, $\varepsilon_{yy}=\gamma_{xy}=0$ on $0\le y\le h$,
with zero gradient energy and the retained reduced tensors, gives the scaling parameter
\begin{equation}
  x=\frac{3\mu^{2}}{C^{p}_{xx}\,h^{2}\,\kappa^{p}_{yy}}
   =\frac{3\,(10^{-6})^{2}}{1.6946\times10^{11}\cdot(10^{-5})^{2}\cdot
     4.9804\times10^{-10}}
   =3.5546\times10^{-4}
  \label{eq:xflexo}
\end{equation}
at $h=\SI{10}{\micro\meter}$. The coefficient 3 follows from
$\int_0^h y^2\,\mathrm{d}y=h^3/3$ for this prescribed edge-axis strain.
A centered prescribed strain $\kappa_c(y-h/2)$ gives coefficient 12 instead.
Neither field is the clamped, gradient-stiffened eigenmode. Equation~\eqref{eq:xflexo}
is therefore an $h^{-2}$ guide with a configuration-dependent prefactor,
not a quantitative benchmark for the coupled FEM result.
For a given fractional order, the L1 weights are determined by
\begin{equation}
  \alpha=0.5:\; (b_0,b_1,b_2)=(1,\;0.414214,\;0.317837),\qquad
  \alpha=0.7:\; (b_0,b_1,b_2)=(1,\;0.231144,\;0.159245),
  \label{eq:bweights}
\end{equation}
and the damping ratio follows from the decrement of two successive peaks,
\begin{equation}
  \delta=\ln\frac{A_k}{A_{k+1}},\qquad
  \zeta_{\rm eff}=\frac{\delta}{\sqrt{4\pi^{2}+\delta^{2}}} ,
  \label{eq:zeta}
\end{equation}
of which $\zeta=\delta/2\pi$ is the small-damping approximation. At
$\omega\tau=0.3$ the ratio of two successive voltage peaks is $2.5947$, so
$\delta=0.95345$ and $\zeta_{\rm eff}=0.15003$ against the exact
$\omega\tau/2=0.15000$.

%% Summary of the worked calculation: inputs -> intermediate -> output.
\begin{table}[!htbp]
\centering
\small
\caption{The headline numbers of the abstract, from their inputs. Each value
is the one carried into the figures and tables of this section; the last column
is the number the abstract quotes.}
\label{tab:worked}
\begin{tabular}{@{}p{0.20\textwidth}p{0.22\textwidth}p{0.24\textwidth}p{0.24\textwidth}@{}}
\toprule
quantity & inputs & intermediate & value \\
\midrule
extensional coupling & $C_p$, $e^{p}$, $\kappa^{p}$ &
$x=1.3542$; $x_{\rm FEM}=1.3402$ & $0.5727$ ($0.5752$ closed form) \\
thickness independence & $k^2$ at $10$ and \SI{1}{\micro\meter} &
$0.57268\to0.57310$ & $0.07\,\%$ \\
flexoelectric growth & $\mu$, $C^{p}_{xx}$, $h$, $\kappa^{p}_{yy}$ &
$1.6677\times10^{-4}\to4.0400\times10^{-3}$ & factor $24$ \\
flexoelectric slope & the eight thicknesses & local log--log slopes &
$-1.96$ to $-0.44$ \\
window anisotropy & $M_{11}$, $M_{22}$, $\theta$ & $M$ rotates, the mesh does
not & $14$--$17\,\%$ \\
two configurations & \eqref{eq:xclosed} in both & $0.4781$ against $0.1036$;
\SI{143.71}{\volt} against \SI{2.22}{\volt} & factor $65$ \\
memory ratio & $K_{\rm eff}$, $K_{uu}$ & $\zeta=0.15003$ against $0.06465$ &
$0.431$ \\
integer-order limit & $\omega\tau=0.3$ & $\zeta_{\rm eff}$ against
$\omega\tau/2$ & $0.02\,\%$ \\
\bottomrule
\end{tabular}
\end{table}


Subsequent calculations retain the specified material, window and mesh settings unless a parameter is explicitly varied.

\FloatBarrier
